\cleardoublepage
\vspace{2cm}
\begin{center}
{\fontsize{14}{16.8}\selectfont\textbf{DAFTAR SIMBOL}}\\[1em]
\end{center}
\label{chap:simbol}
\addcontentsline{toc}{chapter}{DAFTAR SIMBOL}
\setcounter{chapter}{0}
\setcounter{section}{0}
\vspace{2em}

\begin{longtable}{@{} l >{\raggedright\arraybackslash}p{6.5cm} c @{}}
\multicolumn{1}{l}{\textbf{Notasi}} & \multicolumn{1}{l}{\textbf{Deskripsi}} & \multicolumn{1}{l}{\parbox[c]{3.5cm}{\raggedright \textbf{Pemakaian pertama kali}}} \\

% \multicolumn{2}{l}{\textbf{Simbol Polutan dan Kualitas Udara}} \\

PM\textsubscript{10} & Partikulat debu halus dengan ukuran diameter kurang dari 10 mikrometer & 1 \\
CO & Gas Karbon Monoksida & 2 \\
O\textsubscript{3} & Gas Ozon & 2 \\
SO\textsubscript{2} & Gas Sulfur Dioksida & 6 \\
NO\textsubscript{2} & Gas Nitrogen Dioksida & 6 \\[1ex]

% \multicolumn{2}{l}{\textbf{Notasi Waktu dan Deret Waktu}} \\

$t$ & Indeks urutan waktu saat ini (observasi berjalan) & 2 \\
$t+1$ & Proyeksi waktu satu hari ke depan & 2 \\
$t-k$ & Rentang waktu observasi sebelumnya dengan jeda waktu $k$ & 14 \\
$w$ & Rentang jendela observasi berjalan (\textit{rolling window}) & 20 \\
$T$ & Total waktu observasi deret waktu & 15 \\
$Y_t$ & Deret waktu observasi aktual pada waktu $t$ & 14 \\
$X_t$ & Matriks data variabel observasi polutan historis & 36 \\
$L$ & Operator bergeser mundur (\textit{backshift operator}) & 17 \\[1ex]

% \multicolumn{2}{l}{\textbf{Komponen Uji Stasioneritas (ADF)}} \\

$\Delta$ & Operator selisih atau pembedaan (\textit{differencing}) & 14 \\
$\alpha$ & Parameter konstanta pada model regresi uji ADF & 14 \\
$\beta$ & Parameter koefisien tren linier pada uji ADF & 14 \\
$\gamma$ & Koefisien parameter jeda pada uji ADF & 14 \\
$\delta_i$ & Koefisien selisih pada lag ke-$i$ di uji ADF & 14 \\
$\epsilon_t$ & Sisaan acak (\textit{error}) atau \textit{white noise} pada waktu $t$ & 14 \\[1ex]

% \multicolumn{2}{l}{\textbf{Parameter Model Regresi dan Hibrida}} \\

$N$ & Total jumlah baris data observasi & 15 \\
$M$ & Jumlah variabel fitur yang dievaluasi & 15 \\
$B$ & Jumlah ansambel pohon pada algoritma \textit{Random Forest} & 16 \\
$K$ & Jumlah pohon keputusan (\textit{estimators}) & 23 \\
$p$ & Orde komponen \textit{AutoRegressive} (AR) pada ARIMA & 16 \\
$d$ & Orde komponen \textit{Integrated} (I) pada ARIMA & 16 \\
$q$ & Orde komponen \textit{Moving Average} (MA) pada ARIMA & 16 \\
$c$ & Parameter konstanta model ARIMA & 16 \\
$\phi$ & Koefisien pembobot komponen \textit{AutoRegressive} & 16 \\
$\theta$ & Koefisien pembobot komponen \textit{Moving Average} & 16 \\
$N_t$ & Komponen fluktuasi non-linier dari dekomposisi data & 17 \\
$L_t$ & Komponen autokorelasi linier dari dekomposisi data & 17 \\
$e_t$ & Sisa galat prediksi (\textit{residual}) pada waktu $t$ & 18 \\
$y_t'$ & Data hasil pembedaan (\textit{differencing}) tingkat pertama & 16 \\
$\hat{N}_t$ & Estimasi proyeksi nilai awal dari model fase primer & 18 \\
$\hat{L}_t$ & Estimasi taksiran sisaan (\textit{residual}) dari model fase sekunder & 18 \\
$\hat{y}_{RF}$ & Estimasi regresi dari algoritma \textit{Random Forest} & 16 \\
$\hat{r}_{ARIMA}$ & Taksiran sisaan (\textit{residual}) dari komponen statistik ARIMA & 18 \\
$\hat{Y}_{hybrid}$ & Taksiran akhir prediksi dari arsitektur model hibrida sekuensial & 18 \\
$f(\cdot)$ & Fungsi transformasi regresi prediktif & 36 \\[1ex]

% \multicolumn{2}{l}{\textbf{Statistik, Metrik, dan AHP}} \\

$R^2$ & Koefisien determinasi & 2 \\
$\bar{y}$ & Rata-rata nilai observasi data & 15 \\
$\rho_k$ & Koefisien fungsi autokorelasi (ACF) pada jeda waktu $k$ & 15 \\
$MA_w$ & Nilai rata-rata bergerak (\textit{Moving Average}) pada rentang waktu $w$ & 20 \\
$STD_w$ & Nilai standar deviasi bergerak pada rentang waktu $w$ & 20 \\
$RM_w$ & Nilai rata-rata observasi berjalan pada rentang waktu $w$ & 37 \\
$RStd_w$ & Nilai standar deviasi observasi berjalan pada rentang waktu $w$ & 37 \\
$z$ & Skor standar (Z-Score) untuk normalisasi fitur & 54 \\
$\mu$ & Rata-rata (\textit{mean}) distribusi data & 16 \\
$\sigma$ & Deviasi standar distribusi data & 54 \\
$n$ & Jumlah dimensi kriteria matriks AHP / jumlah sampel evaluasi & 41 \\
$\lambda_{max}$ & Nilai Eigen Maksimum pada matriks kriteria AHP & 41 \\
$CI$ & Indeks Konsistensi (\textit{Consistency Index}) pada uji AHP & 41 \\
$RI$ & Indeks Acak Standar (\textit{Random Index}) pada uji AHP & 41 \\
$CR$ & Rasio Konsistensi (\textit{Consistency Ratio}) pada uji AHP & 41 \\[1ex]

% \multicolumn{2}{l}{\textbf{Notasi Lain-lain}} \\

$O(\cdot)$ & Notasi \textit{Big-O} untuk batas atas kompleksitas komputasi & 23 \\
$S$ & Himpunan subset fitur pada perhitungan \textit{Shapley Value} & 24 \\
$\Delta f_i(S)$ & Diferensiasi kontribusi marjinal fitur tunggal (SHAP) & 24 \\
$\phi_i$ & Nilai akhir bobot atribusi SHAP untuk fitur $i$ & 24 \\
$\phi_0$ & Nilai prediksi rata-rata (\textit{base value}) pada algoritma SHAP & 58 \\

\end{longtable}